\subsection{A uniform normalized inverse on finite contours}
\label{sec:finite-linear}

The infinite-end estimates now have to be combined with the long
prescribed paths.  Compact observations remove the Gaussian kernel;
the weighted end estimate must also prevent an approximate kernel
from escaping toward a receding endpoint.  We use the compact norms
and coefficient maps of \Cref{gauge:compact}, and the finite graph
calculations of \Cref{gauge:mixed,gauge:height}.  The end weights are
retained until the final conversion to physical graph bounds.

Here is the geometry behind the parameters in the next proposition.
The large number $B$ measures axial distance along a cylindrical end;
in the later application it is the inverse square root of the opening
scale.  The fixed number $D$ bounds the size of a contour after division
of axial distances by $B$.  A polar height chart writes the surface as
$X=(Zg(r,\theta),r e_r(\theta))$, with $|Z|\asymp B$ and
$e_r(\theta)=(\cos\theta,\sin\theta)$.  Its scalar variable $g$ is
height divided by $Z$, and $r_0$ is its entrance radius; decreasing $r$
corresponds to increasing $\log(r_0/r)$.  A prescribed complex path is
given by $r=r(s)$ on a real outward parameter interval.  The complex
quantity $\log(r_0/r(s))$ occurs only through its coefficients and
chain rule; it is not a new spatial domain.
An outgoing derivative condition prescribes the first derivative of
this scalar graph variable at the terminal section of that interval.
The proposition states the precise analytic hypotheses independently
of any particular contour construction.

\begin{proposition}[Weighted normalized inverse]\label{lin:global}
Fix an angular order $h$, a finite geometric parameter $D$, a cylindrical
weight rate $\alpha_{\mathrm{cyl}}>2.1$, and the
compact observations from Lemma~\ref{lin:observations}.  Consider surfaces of
parameter size $B\to\infty$, obtained by replacing far cylindrical
pieces by real-parameter paths with complex coefficients, optionally
followed by height charts of length $O(\log B)$.  The conical ends may
remain infinite.  The following are the hypotheses on the linear
operators and prescribed charts.

\begin{enumerate}
 \item On each fixed compact subset, the coefficients, compact
 multiplier sources, and observations converge smoothly to
 \eqref{lin:augmented}, or to a member of a sufficiently small uniformly
 invertible perturbation of that block.
 \item From a fixed entrance to $z=cB$, the cylindrical operator is a
 sufficiently small perturbation of
 \eqref{lin:cylinder-operator}.  Choose $c>0$ so that a propagation
 rate $2.1$ is available.  Beyond that point, bounded bands of length
 depending on $D$ satisfy Lemma~\ref{lin:strip}.  A height chart is the
 real-parameter pullback of $\tau=\log(r_0/r)$, with
 $E\asymp(B|r|)^{-2}$, ending with
 $B|r|\ge Y_0$.  Its coefficients obey the same accretivity, derivative,
 lower-term, and small-perturbation bounds.
 \item Each modified cylindrical end has one ordered real outward
 interval $[s_0,T_B]$, equal to the original log coordinate on its
 long real part.  Local real outward coordinates on overlaps are
 $t_i=\chi_i(s)$ with $0<c\le\chi_i'\le C$ and bounded required
 fixed derivatives; the real angular coordinate is unchanged.
 There is a positive scale $E_B(s)$ comparable to the local scales
 with bounded logarithmic derivatives.  The prescribed scalar
 variables are descriptions of one scalar variable through transition
 factors $A_i(s,\theta)$ on fixed overlap bands.  The functions
 $A_i,A_i^{-1}$ and their ordinary real-coordinate derivatives through
 the required fixed orders are bounded.  Choose this scalar variable
 to agree with each prescribed long-chart variable off those bands and
 with the prescribed exit variable near $T_B$.  After these changes,
 the normalized reference operators and their source rows are smooth
 restrictions of one operator on the interval; any additional terms
 from the actual backgrounds satisfy the bounded lower-term and small
 domain-to-range perturbation hypotheses in (ii).  These requirements
 exclude distributional seams.  The coefficients on conical ends satisfy
 Lemma~\ref{lin:conical}, up to a small domain-to-range perturbation.
 \item For each fixed $B$ and sufficiently remote finite stops on
 the conical ends, fixed domain and range isomorphisms identify the
 actual stopped operator with an operator connected to a real
 coercive scalar elliptic problem by an operator-norm continuous homotopy
 on the fixed compact $H^2$ domain with homogeneous boundary
 conditions and range $L^2$. Its coefficients and boundaries are smooth,
 its symbols are properly elliptic, and its boundary conditions are
 complementing throughout. The equal-dimensional
 compact observation and multiplier augmentation is retained.
 This is only an index hypothesis: the uniform estimates in the
 preceding items are required for the actual operators, not for
 the artificial operators along the homotopy.
\end{enumerate}

For clarity, the coefficient requirement in (iii) includes the tangent
of every prescribed complex log path.  If $t=t(s)$ is such a log and
$H=t_s$ is its complex tangent, independent of $\theta$, then the chain
rule on the real $s$ interval gives
\begin{equation}\label{lin:complex-log-pullback}
 \begin{split}
 H\mathcal L={}&\partial_s-\frac aH\partial_{ss}
       -2b\partial_{s\theta}-Hc\partial_{\theta\theta}\\
 &+\left(d_1+\frac{aH_s}{H^2}\right)\partial_s
       +Hd_2\partial_\theta+Hd_0,
 \end{split}
\end{equation}
and the source is multiplied by $H$.  Thus $H,H^{-1}$ and its required
fixed derivatives, including $H_s$, must be bounded.  If the full drift
is divided out, its divisor is
$\delta_H=1+d_1+aH_s/H^2$, and the actual inequalities in (ii) are
\[
 \operatorname{Re}\frac{a}{H\delta_H}\ge cE_B,\quad
 \operatorname{Re}\frac{Hc}{\delta_H}\ge c,\quad
 \left|\frac{a}{H\delta_H}\right|\le CE_B,\quad
 \left|\frac{Hc}{\delta_H}\right|\le C,\quad
 \left|\frac b{\delta_H}\right|\le\varepsilon_*\sqrt{E_B}.
\]
The divided source is $Hf/\delta_H$.  A nonzero divisor alone does not
imply these real-part inequalities.  If $\delta_H=1+O(E_B)$, its
closeness to one preserves previously fixed strict margins.

Choose a weight $\rho_B$ with rate $\alpha_{\mathrm{cyl}}$ on the long real cylinder,
a sufficiently large fixed rate $d_D$ on the bounded modified bands,
and a sufficiently large fixed rate $d_{\mathrm{cap}}$ on a height
chart, retaining its value at the entrance of each new chart.  All
rates exceed their propagation rates by fixed positive margins, are at
least $\alpha_{\mathrm{cyl}}$, and
their interpolations occupy fixed overlap bands where the adjacent
coefficient estimates both hold.  Choose $Y_0$ after these
rates so that $d_{\mathrm{cap}}\sup E$ is sufficiently small; take $B$
large only after these choices.

Here $u$ denotes the scalar unknown in each prescribed chart: normal
height on the core, radius on cylindrical charts, relative height on
cap charts, and link displacement on conical charts.  The corresponding
source uses the same change of graph units and drift normalization.
Let $X_B$ consist of the compact norm \eqref{lin:compact-norm}, the
cylindrical and height norms \eqref{lin:weighted-strip}, the conical
norms \eqref{lin:conical-norm}, and the Euclidean norm of the compact
multiplier vector.  Let $Y_B$ consist of the corresponding compact
and end source norms and the Euclidean observation norm.  Homogeneous
outgoing graph-derivative conditions are included in the domain.
Then
\begin{equation}\label{lin:global-estimate}
 \|(u,\lambda)\|_{X_B}
 \le C_{D,Y_0,h,\alpha_{\mathrm{cyl}}}
 \left\|\bigl(L_Bu+\sum_i\lambda_i\chi_{i,B},P_Bu\bigr)\right\|_{Y_B},
\end{equation}
and this normalized operator is onto.  At a true terminal section define
$g=u_s(T_B)$ for the prescribed scalar variable in the real outward
coordinate.  If $t=t(s)$ is a local log coordinate, this means
$g=H u_t(T_B)$ with $H=t_s$.  For this prescribed exit derivative
the right hand side is augmented by
$\rho_B(T_B)^{-1}\|g\|_{\mathcal T^N_{E(T_B),h}}$ at each terminal section;
the sum over the fixed finite list of exits is understood.
The thresholds and constant may depend on the fixed rate, but are independent
of $B$ and of the lengths of the exhausted
conical ends.  Removing $\rho_B$, or converting a relative height
to physical height on a cap, costs at most a fixed power of $B$.
The preliminary construction uses $\alpha_{\mathrm{cyl}}=2.2$; a source
with a fixed higher Taylor growth rate is permitted by choosing
$\alpha_{\mathrm{cyl}}$ equal to or larger than that rate.  The strict
buffer is above the propagation rate, so equality with the source growth
rate is allowed in the uniformly local source norm.
\end{proposition}

\begin{proof}
We first prove one exterior estimate on the real interval in (iii).
This avoids imposing boundary data at changes of graph gauge.  For a
scalar change $u=A\widetilde u$, conjugate the source by $A^{-1}$.
The principal coefficients of \eqref{lin:admissible-operator} are
unchanged, while the lower coefficients become
\begin{align*}
 \widetilde d_1&=d_1-2aA_s/A-2bA_\theta/A,\\
 \widetilde d_2&=d_2-2bA_s/A-2cA_\theta/A,\\
 \widetilde d_0&=d_0+(1+d_1)A_s/A+d_2A_\theta/A
       -aA_{ss}/A-2bA_{s\theta}/A-cA_{\theta\theta}/A.
\end{align*}
The $a$ contribution to $\widetilde d_1$ is $O(E_B)$; the
$b$ contribution is $O(\varepsilon_*\sqrt{E_B})$ and is an allowed
small strip operator perturbation.  The remaining additions are bounded
lower terms.  In particular $A_s/A$ must have an ordinary bounded
derivative norm, since it occurs without an $E_B$ factor in the mass.
The transition is taken in the prescribed reference frames; the
coefficient--module bounds control the rougher actual backgrounds.
For the one-dimensional real coordinate changes in (iii), the chain
rule is \eqref{lin:complex-log-pullback} with positive real $H$.
Their bounded factors preserve the radial, mixed, and angular units.

The fixed-radius radial-to-height transition illustrates why these
requirements hold for the later circular overlap.  Write
$z=Z_0\varphi(r)$ with $|Z_0|\asymp|Z|\asymp B$ and
$\varphi^2=1-r^2/2$, and put
$H_r=(\log(r_0/r))_s$.  Then
\[
 (\log z)_s=\frac{r^2}{2\varphi^2}H_r,\qquad
 u_{\rm rel}
       =-\frac{z_s}{Zr_s}u_{\rm radius}
       =\frac{Z_0}{Z}\frac r{2\varphi}u_{\rm radius}.
\]
Here $u_{\rm rel}$ is relative height, so the factor follows by equating
the two variations' normal components.
The last factor and its inverse have bounded fixed derivatives when
$|r|\asymp r_0>0$ and $\varphi$ is a unit.  They need not do so at
$|r|\asymp B^{-1}$, so the transition is made on that fixed-radius
overlap.  In the common real parameter, the circular radial and angular
diffusions, before the small axial-curvature correction to the drift, are
\[
 \frac{4}{Z_0^2r^2H_r(1+4\varphi^2/(Z_0^2r^2))},
       \qquad \frac{H_r}{\varphi^2}.
\]
These have the polar units $E_B\asymp(B|r|)^{-2}$ and $1$.
Equivalently, in the mixed gauge of Lemma~\ref{gauge:mixed}, the
angular-curvature and support terms have slow drift
$-\varphi^2/(r^2H_r)$, independent of the prescribed transverse
direction.  The formulas follow from
$r_s=-rH_r$, $z_s=Z_0r^2H_r/(2\varphi)$ and the metric in
\eqref{gauge:metric}.  The remaining axial drift is $O(E_B)$ on this
overlap.  On a long height chart \eqref{lin:complex-log-pullback} and
the strict margins in (ii) are checked with its own scale
$(B|r|)^{-2}$, including the tangent derivatives.  No uniformly bounded
conversion to physical normal height along the whole cap is needed here.

Trivialize on the fixed overlap bands as in (iii), keeping the prescribed
exit variable.  Put $u=\rho_B w$ on the entire interval and
$d=(\log\rho_B)_s$.  On the long real cylinder the model shifted mass at
$d=\alpha_{\mathrm{cyl}}$ is
\[
 d-2-2\mu^2(d^2-d).
\]
It has a strict positive buffer for a far entrance depending on this fixed
rate; at $d=2.2$ the limiting buffer is $0.2$.  The small full
operator-norm perturbation in (ii) preserves this estimate by absorption;
decay of that perturbation is not needed on the interval $z\le cB$.
Raise the rate on a fixed band while this margin still holds, before
entering a region whose lower terms require $d_D$, and similarly enter
the height rate.  On the other bands the algebraic mass is
\[
 (1+\widetilde d_1)d-a(d^2+d_s)+\widetilde d_0.
\]
The chosen rate exceeds the bounded reaction costs, and $E_Bd\ll1$
controls both the quadratic term and the smooth interpolation term.
The coefficient--module perturbations are absorbed with the strict
local margins.  Perform the integration in the proof of
\eqref{lin:forward} once on this whole interval, including the same
additional exponentially decreasing localization centered at an arbitrary
strip.  There are only the true entrance and terminal boundary terms.
The exit variable is unchanged there, so its conjugated condition is
$w_s+d w=\rho_B(T_B)^{-1}g$ and has the favorable terminal sign already
computed.  The geometric strip sum is uniform in the interval length.
Local patch estimates recover every bulk term across the smooth overlap
bands.  We obtain
\begin{equation}\label{lin:exterior-global}
 \|u\|_{X_B(\mathrm{exterior})}
 \le C\left(C_h^2(u;K_{\mathrm{collar}})
       +\|L_Bu\|_{Y_B(\mathrm{exterior})}
       +\|g/\rho_B(T_B)\|_{\mathcal T^N_{E(T_B),h}}\right).
\end{equation}
The overlap constants depend only on the prescribed fixed data and $D$.
Lemma~\ref{lin:conical} supplies the
conical part of this estimate.  The compact multiplier sources are
absent on all sufficiently far collars.

Interior elliptic estimates on nested fixed collars, including the
commutations in \eqref{lin:compact-norm}, now yield
\begin{equation}\label{lin:compact-remainder}
 \|(u,\lambda)\|_{X_B}
 \le C\left(\|\mathcal A_B(u,\lambda)\|_{Y_B}
                    +\|u\|_{L^2(K')}+|\lambda|\right)
\end{equation}
for a fixed compact $K'$.  In particular the estimate has been proved
in the weighted norm, before any power of $B$ is spent.

For a fully stopped problem write $T$ for the smallest logarithmic
stopping coordinate on the conical ends.  If there are no conical
ends, omit $T$.  Suppose there were no thresholds $B_0,T_0$ and
constant $C$ giving \eqref{lin:global-estimate} whenever
$B\ge B_0$ and $T\ge T_0$.  We could then choose a failing sequence
with $B\to\infty$ and every conical stop tending to infinity.
Normalize its unknowns to have $X_B$ norm one
and its output to tend to zero.  Interior compactness and finite
dimensionality of the multipliers give a local limit $(u_\infty,
\lambda_\infty)$.  On every fixed cylindrical interval the weight is
eventually $e^{\alpha_{\mathrm{cyl}}(t-t_0)}$.  Consequently $u_\infty$ has at most
polynomial growth there.  On conical ends its ambient normal height
has at most linear growth, by \eqref{lin:conical-norm}.  The derivative
estimates give $u_\infty\in H^2_G$.  It solves the homogeneous limiting
normalized block and hence vanishes, together with its multipliers.
Local elliptic estimates upgrade this convergence sufficiently to
make the compact terms in \eqref{lin:compact-remainder} vanish.  The
latter inequality contradicts the normalization.  This is the
exclusion of escape along a long end: it uses the fixed strict buffer
in \eqref{lin:exterior-global}, and would not follow by normalizing an
unweighted estimate with a factor $B^q$.
The contradiction gives one bound on the entire quadrant
$B\ge B_0$, $T\ge T_0$.  Thus the subsequent conical exhaustion
may fix any $B\ge B_0$ before sending the conical stops to infinity.

Surjectivity is a separate argument. Stop all ends, including the
conical ends, at remote finite sections. At each fixed $B$ and each
fixed conical stopping radius, the compact scalar boundary problem is
realized from its homogeneous mixed-boundary $H^2$ domain to $L^2$.
The full boundary-data map is Fredholm by the complementing-boundary theorem
\cite[Chapter~1, Section~1, subsection~3, Theorem~1.1]{NP1994}, localized in
surface charts. Entrance Dirichlet and transverse derivative
conditions occur on disjoint boundary components $\Gamma_D,\Gamma_N$.
To pass to the stated
homogeneous realization, let $\mathcal B$ be the boundary trace map into
$\mathcal T=H^{3/2}(\Gamma_D)\oplus H^{1/2}(\Gamma_N)$ and choose its
bounded right inverse $R$ from fixed smooth collars.  The decomposition
$u=v+Rg$, $v\in\ker\mathcal B$, identifies the full map with
$(v,g)\mapsto(L_Bv+L_BRg,g)$.  A bounded triangular change in the range
reduces it to $\operatorname{diag}(L_B|_{\ker\mathcal B},I_{\mathcal T})$,
so the homogeneous realization is Fredholm with the same index.
Hypothesis (iv)
identifies its index with that of a real coercive problem, hence zero.
Continuity of the index and invariance under finite-rank perturbations
are the Banach-space facts in
\cite[Chapter~IV, Sections~5.2--5.3, Theorems~5.17 and~5.26]{Kato1995}.
Adding equally many multiplier variables and observation equations first
gives $\operatorname{ind}\operatorname{diag}(L_B,0_{\mathbb C^m})
=\operatorname{ind}L_B$.  The off-diagonal observation and source maps
are finite rank, so the augmented block still has index zero.

Apply the already proved weighted a priori estimate to the actual
stopped operator. Its constants are independent of the conical
stopping radii by Lemma~\ref{lin:conical} and the compact-limit
argument above. For sufficiently remote stops and large $B$ it
excludes a kernel: elliptic boundary regularity
\cite[Chapter~1, Section~1, subsection~3, Theorem~1.5]{NP1994}
first makes every
kernel element smooth, so the commuted estimate applies. The stopped
augmented operator is therefore bijective on this $H^2$--$L^2$
realization. For smooth sources its solution is smooth; the actual
commuted estimates then bound it in $X_B$ by the source in $Y_B$.
Approximation by smooth sources extends this inverse to the stated
angularly commuted spaces, with the same uniform bound. Thus no
identification of those spaces with an isotropic higher Sobolev
scale is needed. Let the conical stops tend to infinity, keeping $B$
fixed. The uniform conical estimates, weak compactness in their
weighted bulk spaces, and local elliptic compactness pass these
solutions to the infinite-conical domain. The prescribed compact
observations and finite cylindrical or height exits pass to the
limit as well. This proves surjectivity of the actual operator
with the asserted bound.

Only compact stopped problems enter the artificial index homotopy.
No uniform inverse along that homotopy, no geometric realization
of its intermediate operators, and no deformation of an infinite
conical boundary problem are required. Fixed-parameter elliptic
Fredholm theory supplies the index; the weighted estimates for the
actual geometric operators supply uniformity in $B$ and in the
conical exhaustion.

Finally the real-cylinder weight is at most $CB^{\alpha_{\mathrm{cyl}}}$, and bounded bands
cost a fixed factor depending on $D$.  By the stated length hypothesis and
the bounded real pullbacks, a cap has $s$-length at most
$C_{\mathrm{par}}\log B+C_{\mathrm{par}}$, with $C_{\mathrm{par}}$ fixed
before $B$.  Thus
$\sup\rho_B\le C_{D,Y_0,\alpha_{\mathrm{cyl}}}
 B^{\alpha_{\mathrm{cyl}}+C_{\mathrm{par}}d_{\mathrm{cap}}}$.
Conversions of the prescribed height units and any fixed number of
their derivatives add only fixed powers.  This proves the last
assertion without altering the uniform weighted estimate.
\end{proof}

\begin{lemma}[Exponential backward shielding]\label{lin:shielding}
On a fixed log-length overlap suppose the hypotheses of
Lemma~\ref{lin:strip} hold with $E\asymp B^{-2}$.  A solution of the
homogeneous difference equation with zero entrance value and
polynomially bounded farther data is, on every fixed earlier band
separated from those data, bounded by
\begin{equation}\label{lin:shielding-estimate}
 C B^M\exp(-cB^2)
\end{equation}
in the strip norm, with every fixed number of admissible derivatives.
The assertion is stable under the small coefficient perturbations in
the same shifted domain-to-range norms.  It holds with actual outgoing
derivative data as well as with a farther Dirichlet trace.
For holomorphic finite-parameter derivatives, the solution and the
displayed bound are assumed on an outer parameter ball and the derivatives
are evaluated on an inner ball whose distance from the outer boundary is at
least $cB^{-q}$ for fixed $c,q>0$.  For smooth contour derivatives,
the pulled-back differentiated operators and data must separately obey
the corresponding shifted bounds, with a fixed separation of the two
bands.
\end{lemma}

\begin{proof}
Put $Z=B$, $d=c_1Z^2$, $D=\partial_t+d$, and
$v=e^{d(t-T)}w$, with $d$ constant on the overlap.
Choose $c_1>0$ after the reference margins so that $d\sup E$ is
small.  All bulk norms below mean $L^2_tH^h_\theta$ on the fixed
overlap, or on a fixed enlarged patch.  The shifted reference energy
has mass $d-O(Ed^2)\ge cd$; the symbol and boundary argument of
\eqref{lin:forward} control
\[
 N_d(w):=d\|w\|+\|Dw\|+\|w_{\theta\theta}\|
           +\|ED^2w\|+\|\sqrt E D w_\theta\|.
\]
Mixed and imaginary terms use the small derivative energies exactly
as in that argument.  For clarity the first traces also have constants
independent of $Z$.  With $E$ replaced by its comparable patch value,
product integration gives
\begin{align*}
 E\|Dw\|_{L^\infty H^h}^2
 &\le C\bigl((E+Ed)\|Dw\|^2+\|Dw\|\,\|ED^2w\|\bigr),\\
 \|w_\theta\|_{L^\infty H^h}^2
 &\le C\|w_{\theta\theta}\|
                   \bigl((1+d)\|w\|+\|Dw\|\bigr).
\end{align*}
The ordinary product identity supplies the $w$ trace.  Hence
$(w,w_\theta,\sqrt E Dw)$ is bounded by $CN_d(w)$ in
$\mathcal A=L^\infty_tH^h_\theta$.  The term $d\|w\|$ is
retained in this estimate; no second derivative has a trace.

Explicitly, for a perturbation
$\delta aD^2+2\delta bD\partial_\theta+\delta c\partial_\theta^2
+pD+q\partial_\theta+r$, a sufficient coefficient--module norm is
\[
 \mathfrak M=\|\delta a/E\|_{\mathcal A}
 +\|\delta b/\sqrt E\|_{\mathcal A}+\|\delta c\|_{\mathcal A}
 +E^{-1/2}\|p\|+\|q\|+\|r\|.
\]
The last three norms are bulk norms: the preceding traces give the
operator bound $C\mathfrak M N_d(w)$.  Bounded lower coefficients
may also be estimated directly on their bulk slots.

Here is the shifted perturbation check for the geometric difference
operators used below.  In the mixed prescribed gauges of
Lemma~\ref{gauge:mixed}, suppose the two background perturbations
have norm at most $W$ in \eqref{gauge:norm}, and set
$\varepsilon=ZW\ll1$.  Relative to the circular reference,
\eqref{gauge:schur} and \eqref{gauge:slots} give, uniformly along
the averaged segment,
\[
 \|\delta g^{tt}\|_{\mathcal A}\le CE(\varepsilon+\varepsilon^2),
 \quad\|\delta g^{t\theta}\|_{\mathcal A}\le C\sqrt E\varepsilon,
 \quad\|\delta g^{\theta\theta}\|_{\mathcal A}\le C(W+W^2).
\]
Analytic drift division preserves the corresponding normalized bounds.
Applying these coefficients to $D^2w,Dw_\theta,w_{\theta\theta}$
costs at most $C(\varepsilon+\varepsilon^2)N_d(w)$.
Keeping $D$ intact includes the terms with $d^2w$ and $dw_\theta$.

The averaged Jacobian also differentiates the coefficients multiplying
a background second derivative.  Keep that derivative in its original
unshifted $L^2$ slot.  From \eqref{gauge:schur}--\eqref{gauge:slots},
the resulting coefficients of $Dw,w_\theta,w$ have respective
$L^2_tH^h_\theta$ bounds
$CW$, $C(\varepsilon+\varepsilon^2)$, and $CW$.
For example differentiating the circular radial coefficient costs
$CE$, and $\|u_{tt}\|\le CZ^2W$; differentiating its two-tilt
correction costs $CW$, giving $CZ^2W^2$ in the angular slot.
Use $\|Dw\|_{\mathcal A}\le CZN_d(w)$ and the other first traces
to bound all these module products by
$C(\varepsilon+\varepsilon^2)N_d(w)$.
The lower terms are precisely \eqref{gauge:numerators} and
\eqref{gauge:support}: pure speed squares retain $E$, and the
circular angular and support terms are affine in speed.
They obey the same bound.  Thus the full difference operator satisfies
\[
 \|e^{-d(t-T)}(L_*-L_{\rm ref})e^{d(t-T)}w\|_{L^2H^h}
 \le C(\varepsilon+\varepsilon^2)N_d(w).
\]
This is a module estimate, not pointwise control of every lower
coefficient.  In a pure radial chart the variables in
Lemma~\ref{gauge:radial} give the same argument, with any small
circular opening included in the reference.  The reference margins
and this shifted smallness are checked before applying the lemma:
smallness of the opening alone does not assert $ZW\ll1$ for the
correction relative to that reference.

Keep the actual exit condition throughout:
\begin{equation}\label{lin:robin}
 v_t(T)=g\quad\Longleftrightarrow\quad Dw(T)=g.
\end{equation}
For a reference fast root the lift denominator is
$(\lambda_+-d)+d=\lambda_+$, so the natural norm
\eqref{lin:neumann-trace} retains its gain.
For homogeneous data the principal terminal coefficient is
$d\operatorname{Re}a+(1-2d\operatorname{Re}a)/2=1/2$;
$d_1=O(E)$ adds only $O(E)$.
The full perturbation is inverted on this same Robin domain by a
Neumann series once its displayed relative norm is sufficiently small.
No trace of a rough lower coefficient is introduced.
A scalar conversion $v=A(t)\widetilde v$ changes both the interior
operator and the exit to
$(\partial_t+d+A_t/A)\widetilde w=g/A$.
Bounded invertible $A$ with the prescribed derivative bounds gives
equivalent shifted domains.  Returning to ordinary derivatives costs
only fixed powers of $Z$; commuted versions use the corresponding
differentiated coefficient--module bounds.

One may alternatively take the farther actual trace as Dirichlet
data, or cut off the solution near the far side and impose zero data
beyond the cutoff.  The resulting source is supported a fixed distance
to the right of the band being estimated.  The factor
$e^{d(t-T)}$ suppresses its effect there by $e^{-d\,\mathrm{gap}}$.
The shifted estimates have only polynomial losses, which proves
\eqref{lin:shielding-estimate}.  All operations occur in prescribed
charts.  On the inner parameter ball in the statement, Cauchy's estimate
for a derivative of order $j$ costs at most $C_jB^{qj}$, which preserves
the exponential.  A ball's radius by itself gives no estimate at points
arbitrarily close to its boundary.  Smooth contour parameters instead
require differentiation of the equation and its prescribed real-domain
pullback.  When those differentiated coefficients and data have the
stated polynomial bounds and their sources remain on the farther side,
the same shifted estimate applies to each differentiated equation.
\end{proof}

\begin{remark}\label{lin:two-inversions}
The uniform weighted inverse and its polynomial physical version have
different uses.  A refined exterior nonlinear problem must first be
solved in its own mean and angular-oscillation weights, using the
corresponding small Lipschitz constants.  If its entrance agrees with
a preliminary normalized solution, Lemma~\ref{lin:shielding} makes a
splice residual exponentially small.  Proposition~\ref{lin:global}
can then correct that residual even after all fixed polynomial losses.
Applying a coarse $B^q$ inverse directly to the original exterior
power-sized error would require a separate smallness argument and is
not an implication of the proposition.
\end{remark}
